\documentclass[11pt,a4paper]{article}

\usepackage[T1]{fontenc}
\usepackage[utf8]{inputenc}
\usepackage{lmodern}
\usepackage[a4paper,top=2.4cm,bottom=2.2cm,left=2.2cm,right=2.2cm]{geometry}
\usepackage{amsmath}
\usepackage{amssymb}
\usepackage{graphicx}
\graphicspath{{./}}
\usepackage{array}
\usepackage{makecell}
\renewcommand{\arraystretch}{1.25}
\usepackage{enumitem}
\usepackage{fancyhdr}
\usepackage{xcolor}

\newcommand{\ohms}{\ensuremath{\Omega}}
\newcommand{\degC}{\ensuremath{^{\circ}}C}
\newcommand{\dg}{\ensuremath{^{\circ}}}
\newcommand{\pow}[1]{\ensuremath{^{#1}}}

\newcommand{\topiclabel}[1]{{\hfill\normalfont\small\itshape\textcolor{black!55}{#1}}}

\newlength{\anslineskip}
\setlength{\anslineskip}{8mm}

\newcounter{ansln}
\newcommand{\Alines}[2]{%
  \par\vspace{1mm}%
  \setcounter{ansln}{1}%
  \loop\ifnum\value{ansln}<#1\relax
    \noindent\mbox{}\dotfill\par\vspace{\anslineskip}%
    \stepcounter{ansln}%
  \repeat
  \noindent\mbox{}\dotfill\hspace{0.6em}\mbox{#2}\par
  \vspace{1mm}%
}

\newcommand{\plainline}{\par\noindent\mbox{}\dotfill\par\vspace{\anslineskip}}

\newcommand{\labelline}[1]{\par\vspace{1mm}\noindent #1\hspace{0.5em}\dotfill\par\vspace{\anslineskip}}

\newcommand{\ansval}[3]{%
  \par\vspace{5mm}%
  \noindent\hspace*{3.5cm}#1\hspace{0.5em}\dotfill\hspace{0.5em}#2\hspace{0.5em}\makebox[1.1cm][r]{#3}\par
  \vspace{3mm}%
}

\renewcommand{\markright}[1]{\par\nopagebreak\noindent\hspace*{\fill}\mbox{#1}\par\vspace{1mm}}

% Per-question head: \examq{paper-id}{number}{topics}{marks}
% Renders: paper-id - Q<number> - topics - marks  (one line, same font).
% Same command and same rendered line as the answer paper.
\newcommand{\swmarks}{}
\newcommand{\examq}[4]{%
  \gdef\swmarks{#4}%
  \par\addvspace{16pt}\noindent
  {\large\bfseries #1 - Q#2 - #3 - #4}\par\vspace{5pt}%
}
% \total now takes NO argument: it prints the marks stored by \examq.
\newcommand{\total}{\par\vspace{2mm}\noindent\hspace*{\fill}[Total:\ \swmarks]\par\vspace{2mm}}

\newcommand{\qfig}[3][0.5\linewidth]{%
  \par\vspace{2mm}
  \begin{center}%
    \IfFileExists{#2}%
      {\includegraphics[width=#1]{#2}}%
      {\fbox{\parbox[c][26mm][c]{#1}{\centering\ttfamily\footnotesize #2}}}\\[4pt]
    \textbf{#3}%
  \end{center}%
  \vspace{2mm}\par
}

\newlist{parts}{enumerate}{1}
\setlist[parts]{label=\textbf{(\alph*)},leftmargin=2.1em,labelsep=0.6em,itemsep=10pt,topsep=6pt,parsep=4pt}

\newlist{subparts}{enumerate}{1}
\setlist[subparts]{label=\textbf{(\roman*)},leftmargin=2.3em,labelsep=0.6em,itemsep=10pt,topsep=6pt,parsep=4pt}

\pagestyle{fancy}
\fancyhf{}
\fancyhead[L]{\footnotesize 4024/12/M/J/25}
\fancyhead[C]{\footnotesize Cambridge O Level Mathematics (Syllabus D) -- Paper 1}
\fancyhead[R]{\footnotesize May/June 2025}
\fancyfoot[C]{\footnotesize\thepage}
\renewcommand{\headrulewidth}{0.4pt}
\renewcommand{\footrulewidth}{0pt}

\setlength{\parindent}{0pt}

\begin{document}

\examq{4024/12/M/J/25}{1}{Number Fundamentals \textperiodcentered\ Fractions Fundamentals}{4}
Work out.
\begin{parts}
  \item $6-2\times(-4)$
  \ansval{}{}{[1]}

  \item $80\times0.2$
  \ansval{}{}{[1]}

  \item $\dfrac{2}{9}\div\dfrac{5}{6}$
  \vspace{2cm}
  \ansval{}{}{[2]}
\end{parts}
\total

\examq{4024/12/M/J/25}{2}{Probability Fundamentals}{2}
A bag contains 11 balls.\\
There are 5 blue balls and 4 yellow balls.\\
The rest of the balls are green.\\
A ball is taken from the bag at random.

Find the probability that the ball is
\begin{parts}
  \item yellow
  \ansval{}{}{[1]}

  \item \textbf{not} blue.
  \ansval{}{}{[1]}
\end{parts}
\total

\examq{4024/12/M/J/25}{3}{Geometric Fundamentals}{3}
The area of one face of a cube is $9\,\mathrm{cm}^{2}$.

On the 1\,cm grid, draw an accurate net of the cube.
\qfig[0.92\linewidth]{4024_12_M_J_25_fig1.png}{}
\markright{[3]}
\total
\clearpage
\examq{4024/12/M/J/25}{4}{Angles in Polygons \& Parallel Lines}{3}
\qfig[0.72\linewidth]{4024_12_M_J_25_fig2.png}{}
$AEC$ and $BED$ are straight lines.\\
$ED = EC$.

Find the value of $x$.
\vspace{3cm}
\ansval{$x =$}{}{[3]}
\total

\examq{4024/12/M/J/25}{5}{Linear Equations \& Inequalities}{2}
Solve.
\[5(4-x) = 35\]
\vspace{2.5cm}
\ansval{$x =$}{}{[2]}
\total

\examq{4024/12/M/J/25}{6}{Bearings, Constructions \& Scale Drawings}{4}
The scale drawing shows the positions of two villages, $A$ and $B$.\\
The scale is 1\,cm to 5\,km.
\qfig[0.6\linewidth]{4024_12_M_J_25_fig3.png}{Scale: 1\,cm to 5\,km}
\begin{parts}
  \item Find the actual distance between village $A$ and village $B$.
  \vspace{1.5cm}
  \ansval{}{km}{[2]}

  \item Village $C$ is on a bearing of $060\dg$ from village $A$.\\
  Village $C$ is on a bearing of $320\dg$ from village $B$.

  Find and label the position of village $C$ on the scale drawing.
  \markright{[2]}
\end{parts}
\total

\examq{4024/12/M/J/25}{7}{Powers, Roots \& Standard Form}{3}
Evaluate.
\begin{parts}
  \item $\sqrt[3]{125}$
  \ansval{}{}{[1]}

  \item $4^{-2}$
  \vspace{1cm}
  \ansval{}{}{[2]}
\end{parts}
\total

\examq{4024/12/M/J/25}{8}{Scatter Graphs \& Correlation}{4}
Asha records the distance she walks and the time she takes for each of 10 walks.\\
The table shows her results.

\begin{center}
\small
\begin{tabular}{|l|c|c|c|c|c|c|c|c|c|c|}
\hline
Distance (km) & 4.5 & 4.6 & 7.2 & 8.4 & 5.5 & 7.5 & 4.2 & 9.0 & 3.8 & 5.6 \\
\hline
Time (minutes) & 52 & 60 & 93 & 105 & 65 & 100 & 52 & 116 & 49 & 62 \\
\hline
\end{tabular}
\end{center}

\begin{parts}
  \item \leavevmode
  \qfig[0.72\linewidth]{4024_12_M_J_25_fig4.png}{}
  Complete the scatter diagram.\\
  The first 6 points have been plotted for you.
  \markright{[2]}

  \item Draw a line of best fit.
  \markright{[1]}

  \item Asha goes for another walk.\\
  She walks a distance of 6.8\,km.

  Use your line of best fit to estimate the time Asha takes for this walk.
  \ansval{}{minutes}{[1]}
\end{parts}
\total

\examq{4024/12/M/J/25}{9}{Rounding, Estimation \& Bounds}{2}
The diagram shows a rectangle.
\qfig[0.6\linewidth]{4024_12_M_J_25_fig5.png}{}
By writing each number correct to 1 significant figure, find an estimate for the area of the rectangle.
\vspace{3cm}
\ansval{}{mm\pow{2}}{[2]}
\total

\examq{4024/12/M/J/25}{10}{Prime Factors, HCF \& LCM}{3}
\begin{parts}
  \item Write 228 as a product of its prime factors.
  \vspace{4cm}
  \ansval{}{}{[2]}

  \item $228^{2} = 51\,984$

  Write $51\,984$ as a product of its prime factors.
  \vspace{1.5cm}
  \ansval{}{}{[1]}
\end{parts}
\total

\examq{4024/12/M/J/25}{11}{Forming \& Solving Equations \textperiodcentered\ Simultaneous Equations}{5}
The mass of a small box is $x$\,kg.\\
The mass of a large box is $y$\,kg.
\begin{parts}
  \item The total mass of 4 small boxes and 6 large boxes is 30\,kg.

  Show that $\;2x+3y = 15$.
  \vspace{2.5cm}
  \markright{[1]}

  \item The total mass of 6 small boxes and 1 large box is 13\,kg.

  Use this information to write down an equation in terms of $x$ and $y$.
  \vspace{1cm}
  \ansval{}{}{[1]}

  \item Solve the simultaneous equations to find the mass of a small box and the mass of a large box.\\
  You must show all your working.
  \vspace{7cm}
  \ansval{Small box $=$}{kg}{}
  \ansval{Large box $=$}{kg}{}
  \markright{[3]}
\end{parts}
\total

\examq{4024/12/M/J/25}{12}{Transformations}{8}
Triangle $A$ and triangle $B$ are drawn on the grid.
\qfig[0.85\linewidth]{4024_12_M_J_25_fig6.png}{}
\begin{parts}
  \item Describe fully the \textbf{single} transformation that maps triangle $A$ onto triangle $B$.
  \Alines{2}{[3]}

  \item Draw the image of triangle $A$ after a reflection in the line $x = -1$.
  \markright{[2]}

  \item Triangle $A$ is enlarged with scale factor $k$ and centre $(0,\,0)$.\\
  Triangle $C$ is the image of triangle $A$ after the enlargement.\\
  The coordinates of one vertex of triangle $C$ are $(12,\,3)$.
  \begin{subparts}
    \item Find the value of $k$.
    \ansval{$k =$}{}{[1]}

    \item Find the coordinates of the other two vertices of triangle $C$.
    \par\vspace{5mm}\noindent\hspace*{\fill}(\makebox[2.2cm]{\dotfill} , \makebox[2.2cm]{\dotfill}) and (\makebox[2.2cm]{\dotfill} , \makebox[2.2cm]{\dotfill})\hspace{0.6em}\mbox{[2]}\par\vspace{3mm}
  \end{subparts}
\end{parts}
\total

\examq{4024/12/M/J/25}{13}{Percentages Fundamentals}{4}
In a sale, a shop reduces all prices by 20\%.
\begin{parts}
  \item A coat costs \$85 before the sale.

  Work out the sale price of the coat.
  \vspace{4cm}
  \ansval{\$}{}{[2]}

  \item The sale price of a shirt is \$40.

  Work out the cost of the shirt before the sale.
  \vspace{4cm}
  \ansval{\$}{}{[2]}
\end{parts}
\total

\examq{4024/12/M/J/25}{14}{Vectors}{9}
\qfig[0.72\linewidth]{4024_12_M_J_25_fig7.png}{}
$A$ is the point $(-4,\,5)$ and $B$ is the point $(6,\,1)$.
\[\overrightarrow{BC} = \begin{pmatrix} -3 \\ -4 \end{pmatrix}\]
\begin{parts}
  \item Find the column vector $\overrightarrow{AB}$\,.
  \vspace{3cm}
  \par\vspace{5mm}\noindent\hspace*{\fill}$\overrightarrow{AB} = \left(\rule{0pt}{10mm}\hspace{1.5cm}\right)$\hspace{0.6em}\mbox{[2]}\par\vspace{3mm}

  \item Find the coordinates of point $C$.
  \vspace{3cm}
  \par\vspace{5mm}\noindent\hspace*{\fill}(\makebox[2.4cm]{\dotfill} , \makebox[2.4cm]{\dotfill})\hspace{0.6em}\mbox{[2]}\par\vspace{3mm}

  \item $ABCD$ is a trapezium.\\
  $AB$ is parallel to $DC$ and $AB = 2DC$.
  \begin{subparts}
    \item Find the coordinates of point $D$.
    \vspace{3cm}
    \par\vspace{5mm}\noindent\hspace*{\fill}(\makebox[2.4cm]{\dotfill} , \makebox[2.4cm]{\dotfill})\hspace{0.6em}\mbox{[2]}\par\vspace{3mm}

    \item Find the length of line $AD$.\\
    Give your answer as a surd in its simplest form.
    \vspace{5cm}
    \ansval{}{}{[3]}
  \end{subparts}
\end{parts}
\total

\examq{4024/12/M/J/25}{15}{Surds}{3}
\begin{parts}
  \item Simplify.
  \[\sqrt{175}-\sqrt{28}\]
  \vspace{3cm}
  \ansval{}{}{[2]}

  \item Rationalise the denominator.
  \[\dfrac{1}{\sqrt{5}}\]
  \vspace{2cm}
  \ansval{}{}{[1]}
\end{parts}
\total

\examq{4024/12/M/J/25}{16}{Cumulative Frequency}{6}
A group of 80 people each record their journey time from home to work one day.\\
The cumulative frequency diagram shows the results.
\qfig[0.8\linewidth]{4024_12_M_J_25_fig8.png}{}
\begin{parts}
  \item Use the cumulative frequency diagram to find an estimate of
  \begin{subparts}
    \item the median
    \ansval{}{minutes}{[1]}

    \item the interquartile range
    \vspace{1.5cm}
    \ansval{}{minutes}{[2]}

    \item the number of people who had a journey time of 40 minutes or more.
    \vspace{1.5cm}
    \ansval{}{}{[2]}
  \end{subparts}

  \item Each of the 80 people also record their journey time from work to home that day.\\
  The table shows the results.

  \begin{center}
  \begin{tabular}{|l|c|}
  \hline
  Median & 35 minutes \\
  \hline
  Interquartile range & 15 minutes \\
  \hline
  \end{tabular}
  \end{center}

  Jay says:

  ``The journey times from home to work are more consistent than the journey times from work to home.''

  Is Jay correct?\\
  Explain how you decide.

  \par\vspace{3mm}\noindent\makebox[3cm]{\dotfill}\hspace{0.5em}because\hspace{0.5em}\dotfill\hspace{0.6em}\mbox{[1]}\par\vspace{2mm}
\end{parts}
\total

\examq{4024/12/M/J/25}{17}{Rearranging Formulas \textperiodcentered\ Linear Graphs}{5}
The equation of line $L$ is $\;5y+3x = 10\;$.
\begin{parts}
  \item Rearrange $\;5y+3x = 10\;$ to make $y$ the subject.
  \vspace{3.5cm}
  \ansval{$y =$}{}{[2]}

  \item Line $P$ is perpendicular to line $L$.\\
  Line $P$ passes through the point $(6,\,7)$.

  Find the equation of line $P$.
  \vspace{6cm}
  \ansval{}{}{[3]}
\end{parts}
\total

\examq{4024/12/M/J/25}{18}{Real-Life Graphs}{5}
The diagram shows the speed--time graph for a cyclist's journey.
\qfig[0.75\linewidth]{4024_12_M_J_25_fig9.png}{}
\begin{parts}
  \item Describe the motion of the cyclist between $t = 40$ and $t = 100$.
  \Alines{1}{[1]}

  \item The acceleration of the cyclist between $t = 0$ and $t = 40$ is $0.25\,\mathrm{m/s}^{2}$.

  Show that $v = 10$.
  \vspace{2.5cm}
  \markright{[1]}

  \item The total distance travelled by the cyclist between $t = 0$ and $t = T$ is 1.4\,km.

  Find the value of $T$.
  \vspace{6cm}
  \ansval{$T =$}{}{[3]}
\end{parts}
\total

\examq{4024/12/M/J/25}{19}{Expanding \& Factorising Brackets \textperiodcentered\ Algebraic Fractions}{4}
Simplify.
\[\dfrac{3x^{2}-12}{2x^{2}+11x+14}\]
\vspace{6cm}
\ansval{}{}{[4]}
\total

\examq{4024/12/M/J/25}{20}{Fractions, Decimals \& Percentages}{4}
Work out.
\[0.\dot{1}\dot{7}+\dfrac{5}{9}\]
Give your answer as a fraction in its simplest form.
\vspace{7cm}
\ansval{}{}{[4]}
\total

\examq{4024/12/M/J/25}{21}{Vectors}{5}
\qfig[0.85\linewidth]{4024_12_M_J_25_fig10.png}{}
$\overrightarrow{OA} = \mathbf{a}$ and $\overrightarrow{OB} = \mathbf{b}$.\\
$X$ is a point on $AB$ where $AX : XB = 3 : 2$.\\
$OXC$ is a straight line.\\
$AC$ is parallel to $OB$.
\begin{parts}
  \item Find $\overrightarrow{AX}$.\\
  Give your answer in its simplest form in terms of $\mathbf{a}$ and $\mathbf{b}$.
  \vspace{3.5cm}
  \ansval{$\overrightarrow{AX} =$}{}{[2]}

  \item Find $\overrightarrow{XC}$.\\
  Give your answer in its simplest form in terms of $\mathbf{a}$ and $\mathbf{b}$.
  \vspace{4.5cm}
  \ansval{$\overrightarrow{XC} =$}{}{[3]}
\end{parts}
\total

\examq{4024/12/M/J/25}{22}{Quadratic Equations \textperiodcentered\ Quadratic Graphs}{7}
\begin{parts}
  \item Write $\;x^{2}+4x-12\;$ in the form $\;\left(x+a\right)^{2}+b\;$.
  \vspace{3cm}
  \ansval{}{}{[2]}

  \item Use your answer to \textbf{part (a)} to find the coordinates of the turning point of the graph of\\
  $y = x^{2}+4x-12\;$.
  \vspace{1.5cm}
  \par\vspace{5mm}\noindent\hspace*{\fill}(\makebox[2.4cm]{\dotfill} , \makebox[2.4cm]{\dotfill})\hspace{0.6em}\mbox{[1]}\par\vspace{3mm}

  \item Sketch the graph of $\;y = x^{2}+4x-12\;$ indicating the values where the graph crosses the axes.
  \qfig[0.6\linewidth]{4024_12_M_J_25_fig11.png}{}
  \markright{[4]}
\end{parts}
\total

\examq{4024/12/M/J/25}{23}{Circles, Arcs \& Sectors}{5}
\qfig[0.7\linewidth]{4024_12_M_J_25_fig12.png}{}
$OAB$ is a minor sector of a circle, centre $O$.\\
$OCD$ is a major sector of a different circle, centre $O$.

$OCA$ and $ODB$ are straight lines.\\
$OC = 6$\,cm and $OA = 9$\,cm.\\
The length of the minor arc $AB$ is $5\pi$\,cm.

Work out the area of the major sector $OCD$.\\
Give your answer in terms of $\pi$.
\vspace{9cm}
\ansval{}{cm\pow{2}}{[5]}
\total

\end{document}